\documentclass[11pt]{article}
\usepackage[a4paper,margin=2.25cm]{geometry}
\usepackage{amsmath,amssymb}
\usepackage{graphicx}
\usepackage{booktabs}
\usepackage{microtype}
\usepackage{hyperref}
\usepackage{caption}
\usepackage{newtxtext,newtxmath}
\hypersetup{colorlinks=true,linkcolor=black,citecolor=black,urlcolor=blue}
\captionsetup{font=small,labelfont=bf}
\newcommand{\Ti}{$^{49}$Ti}
\newcommand{\DeltaInter}{\Delta_{\mathrm{inter}}^{\min}}
\newcommand{\fR}{f_R}
\title{Supplementary Information\\Quadrupole-limited control windows for an atomically addressable $^{49}$Ti nuclear-spin qudit}
\author{Aishwarya Vishwakarma\\Department of Quantum Matter Physics, University of Geneva, 1211 Geneva, Switzerland\\\href{mailto:Aishwarya.Vishwakarma@unige.ch}{Aishwarya.Vishwakarma@unige.ch}}
\date{}
\begin{document}
\maketitle

\section{Hamiltonian convention and parameter provenance}
The full calculation uses $S=1/2$, $I=7/2$ and
\begin{align}
\frac{H}{h}={}&\frac{\mu_B}{h}\left[g_z(B+B_{\rm tip}\cos\phi)S_z+g_\perp B_{\rm tip}\sin\phi S_x\right]+\frac{\mu_N}{h}g_NBI_z\\
&+A_zS_zI_z+A_\perp(S_xI_x+S_yI_y)+\frac{\kappa}{2I(2I-1)}\left[I_z^2-\frac{1}{2}(I_x^2+I_y^2)\right].
\end{align}
The fitted STM tip field is applied to the electron Zeeman term and not the nuclear Zeeman term, following the physical convention of the single-atom NMR model \cite{vennema2025}. The $^{47}$Ti tip parameters are contextual sensitivity inputs rather than measurements on the \Ti\ readout junction.

\begin{table}[h]
\centering
\caption{Numerical inputs and their role.}
\begin{tabular}{lll}
\toprule
Quantity & Value or range & Role\\
\midrule
$I(^{49}{\rm Ti})$ & $7/2$ & isotope property\\
$A_z$ baseline & 130 MHz & \Ti\ hyperfine scale \cite{stolte2025}\\
$A_z$ comparison & 132.1 MHz & $^{47}$Ti fit \cite{vennema2025}\\
$A_\perp/A_z$ & 0, 0.05, 0.10, 0.20 & sensitivity grid\\
$\kappa_{47}$ & $-56.7\pm0.8$ MHz & $^{47}$Ti fit \cite{vennema2025}\\
$\kappa_{49}$ & $-46.3738\pm2.6554$ MHz & same-EFG isotope scaling\\
$|g_N(^{49}{\rm Ti})|$ & 0.314686 & numerical nuclear factor\\
$g_z$ & $0.56\pm0.02$ & contextual sensitivity range\\
$B_{\rm tip}$ & $67.9\pm2.1$ mT & contextual sensitivity range\\
$\phi$ & $5\pm2^\circ$ & contextual sensitivity range\\
$g_\perp/g_z$ & 0.5, 1.0, 1.5 & sensitivity grid\\
\bottomrule
\end{tabular}
\end{table}

\subsection{Isotope-scaling arithmetic}
Stone \cite{stone2016} gives the recommended ground-state moments
\begin{equation}
Q_{47}=0.302(10)\ {\rm b},\qquad Q_{49}=0.247(11)\ {\rm b}.
\end{equation}
For the same EFG,
\begin{equation}
R_Q=\frac{Q_{49}}{Q_{47}}=0.8178808,
\end{equation}
with relative standard uncertainty
\begin{equation}
\frac{\sigma_R}{R_Q}=\sqrt{\left(\frac{0.011}{0.247}\right)^2+\left(\frac{0.010}{0.302}\right)^2}=0.05550.
\end{equation}
Thus $\kappa_{49}=R_Q\kappa_{47}=-46.37384$ MHz. Combining the $\kappa_{47}$ fit uncertainty and the quadrupole-moment uncertainty in quadrature gives
\begin{equation}
\sigma_{\kappa_{49}}=\sqrt{(R_Q\,0.8)^2+(56.7\,R_Q\,0.05550)^2}=2.6554\ {\rm MHz}.
\end{equation}
The dominant uncertainty in this isotope scaling comes from the nuclear quadrupole moments, not the fitted $\kappa_{47}$ alone.

\section{Explicit derivation of the rigid-comb result}
Using $I_x^2+I_y^2=I(I+1)-I_z^2$, the diagonal axial quadrupole energy is
\begin{equation}
\frac{E_Q(m_I)}{h}=\frac{\kappa}{4I(2I-1)}\left[3m_I^2-I(I+1)\right].
\end{equation}
The signed adjacent transition $m_I\rightarrow m_I-1$ in electron manifold $m_s$ is therefore
\begin{equation}
\nu_{m_s,m_I}=\gamma_NB+A_zm_s+\frac{3\kappa}{4I(2I-1)}(2m_I-1).
\end{equation}
The physical spectroscopy measures $|\nu|$. In the strong-hyperfine regime $|A_z|/2>|\gamma_NB|+3S$, all seven $m_s=+1/2$ transitions retain the sign of $A_z/2$ and all seven $m_s=-1/2$ transitions have the opposite signed orientation. Taking the absolute value therefore reverses the order of the latter comb. With $n,m=-3,\ldots,3$,
\begin{align}
f_n^{(+)}&=\frac{|A_z|}{2}+\gamma_NB+nS,\\
f_m^{(-)}&=\frac{|A_z|}{2}-\gamma_NB-mS,\\
S&=\frac{3|\kappa|}{2I(2I-1)}.
\end{align}
Hence
\begin{equation}
f_n^{(+)}-f_m^{(-)}=2\gamma_NB+(n+m)S
\end{equation}
and
\begin{equation}
\DeltaInter(B)=\min_{\ell}|\ell S-2\gamma_NB|.
\end{equation}
This derivation makes explicit why the common $|A_z|/2$ offset cancels from the observable inter-manifold line separation.

The nearest point of one rigid comb to the other can never lie farther than half a comb spacing, so $\DeltaInter\leq S/2$. Since the nearest intra-manifold spacing is $S$, the inter-manifold constraint is always the tighter one for dual-manifold control.

\section{General spin dependence}
For the axial convention used above,
\begin{equation}
S_I=\frac{3|\kappa|}{2I(2I-1)},\qquad
\Delta_{\max}=\frac{3|\kappa|}{4I(2I-1)},
\end{equation}
and
\begin{equation}
B_r^{\max}=\frac{3(2r+1)|\kappa|}{8I(2I-1)|\gamma_N|}.
\end{equation}
The following dimensionless coefficients illustrate the spin dependence without importing host-specific quadrupole couplings.
\begin{table}[h]
\centering
\caption{Axial comb coefficients for representative high-spin nuclei. The formulas require the host-specific $\kappa$ and $\gamma_N$ for numerical fields.}
\begin{tabular}{lccc}
\toprule
Representative nucleus & $I$ & $S_I/|\kappa|$ & $\Delta_{\max}/|\kappa|$\\
\midrule
$^{173}$Yb & $5/2$ & $3/20$ & $3/40$\\
$^{49}$Ti, $^{123}$Sb & $7/2$ & $1/14$ & $1/28$\\
$^{209}$Bi & $9/2$ & $1/24$ & $1/48$\\
\bottomrule
\end{tabular}
\end{table}
The central point is that the field prescription is not specific to Ti. Once $I$, $\gamma_N$ and the host-dependent $\kappa$ are known, the same comb construction gives the crossing and maximum fields.

\section{Full-model sensitivity envelope and validation}
The sensitivity grid contains $3\times3\times3\times3\times2\times4\times3=1944$ scenarios. It is deterministic and should not be interpreted as a statistical confidence set. The reconstructed implementation was validated against the previous archived values before the low-field extension was accepted:
\begin{center}
\begin{tabular}{lccc}
\toprule
Field & reconstructed worst (MHz) & archived worst (MHz) & difference (MHz)\\
\midrule
1.35 T & 0.0541847 & 0.0541960 & $-1.13\times10^{-5}$\\
1.714 T & 1.0927570 & 1.0927410 & $+1.60\times10^{-5}$\\
\bottomrule
\end{tabular}
\end{center}
The low-field search used 1 mT spacing from 0.95 to 1.11 T, followed by 0.2 mT refinement around the coarse maximum. The result is
\begin{equation}
B_{\rm minimax}=1.0292\ {\rm T},\quad
\min_{\rm grid}\DeltaInter=1.187772\ {\rm MHz},
\end{equation}
with a contiguous coarse-grid interval 0.992 to 1.069 T for which the worst sampled separation is at least 1 MHz.

\subsection{One-at-a-time sensitivity decomposition}
Each documented parameter range was varied independently around the source-backed realization at 1.0292 T. Figure~\ref{fig:oat} and Table~\ref{tab:oat} show the resulting span of $\DeltaInter$.
\begin{figure}[h]
\centering
\includegraphics[width=0.9\linewidth]{figures/FigS1_oat_sensitivity.pdf}
\caption{One-at-a-time span of the full-model source-backed separation at 1.0292 T. The quadrupole coupling and transverse hyperfine ratio dominate the sensitivity by more than two orders of magnitude over the tip-field and angular axes.}
\label{fig:oat}
\end{figure}
\begin{table}[h]
\centering
\caption{One-at-a-time spans at 1.0292 T.}
\label{tab:oat}
\begin{tabular}{lr}
\toprule
Parameter axis & span in $\DeltaInter$ (MHz)\\
\midrule
$\kappa$ & 0.29969\\
$A_\perp/A_z$ & 0.23041\\
$g_z$ & 0.00238\\
$A_z$ & 0.00109\\
$g_\perp/g_z$ & 0.000277\\
$\phi$ & 0.000209\\
$B_{\rm tip}$ & 0.000143\\
\bottomrule
\end{tabular}
\end{table}
The full minimax corner combines several extremes, but the one-at-a-time decomposition shows that the model leverage is concentrated in $\kappa$ and $A_\perp/A_z$.

\section{Fourteen-tone control construction and numerical integration}
For a static electron manifold, the projected nuclear basis has seven adjacent links. The physical waveform contains the seven frequencies from each of the two electron manifolds. A common physical amplitude scale $\fR$ is applied to all tones. For transition $m$, the interaction-picture coupling is proportional to
\begin{equation}
\fR d_m,\qquad d_m=\langle m-1|I_x|m\rangle,
\end{equation}
so the Clebsch-Gordan amplitudes needed for an $I_y$ rotation arise automatically from the projected $I_x$ matrix elements. Each tone phase is $\arg(d_m)+\pi/2$. This is the multitone SU(2) construction used in high-spin nuclear control \cite{asaad2020,fernandez2024,yu2025}.

The target is $U_{\rm target}=\exp[-i(\pi/2)I_y]$ and $t_g=1/(4\fR)$. The interaction-picture $8\times8$ propagator is integrated independently in both electron manifolds under the same waveform using DOP853 with relative tolerance $2\times10^{-8}$ and absolute tolerance $2\times10^{-10}$. The maximum time step is restricted by the gate time and by the largest detuning below 15 MHz. For $d=8$, the average gate fidelity is
\begin{equation}
F_{\rm avg}=\frac{\mathrm{Tr}(U^\dagger U)+|\mathrm{Tr}(U_{\rm target}^\dagger U)|^2}{d(d+1)}
\end{equation}
\cite{nielsen2002}. The reported control metric is the minimum of the two electron-manifold fidelities.

The source-backed realization is $A_z=130$ MHz, $\kappa=-46.3738411$ MHz, $g_z=0.56$, $B_{\rm tip}=67.9$ mT, $\phi=5^\circ$, $A_\perp/A_z=10/130$ and $g_\perp/g_z=1$.
\begin{table}[h]
\centering
\caption{Representative source-backed worst-manifold fidelities.}
\begin{tabular}{rccc}
\toprule
$\fR$ (kHz) & 1.35 T & 1.0292 T & 1.714 T\\
\midrule
5 & 0.988043 & 0.999647 & 0.999475\\
10 & 0.941427 & 0.997627 & 0.997439\\
20 & 0.545803 & 0.992079 & 0.990802\\
35 & 0.198257 & 0.967298 & 0.976625\\
\bottomrule
\end{tabular}
\end{table}
Dense scans use 0.25 kHz spacing around the first 99\% crossing. The interpolated thresholds are 4.8150, 21.9289 and 21.6133 kHz at 1.35, 1.0292 and 1.714 T, respectively. The non-monotonicity close to the threshold is visible in Fig. 2 of the main text. These coherent revivals motivate the first-crossing definition rather than the largest isolated sampled passing point.

A previous diagnostic normalized the two optimized curves by the matched source-backed $\DeltaInter$. Because those two separations differ by less than 1\%, this comparison does not provide a sufficient lever arm to establish a scaling law. It is retained in the archived exploratory analysis but is not used as a publication-facing claim. A genuine $\fR^{99}$ versus $\DeltaInter$ scaling test requires additional control propagations at intermediate fields.

\section{Electron occupation and switching specification}
The single-shot experiment was performed at $T=0.4$ K and reports significant thermal population of the excited electron state \cite{stolte2025}. Using only the external-field source-backed Zeeman splitting at 1.0292 T,
\begin{equation}
\frac{g_z\mu_BB}{h}=8.07\ {\rm GHz},\qquad \frac{h\nu}{k_B}=0.387\ {\rm K},
\end{equation}
which gives the two-level Boltzmann estimate
\begin{equation}
p_{\rm exc}=\frac{1}{1+e^{0.387/0.4}}\simeq0.275.
\end{equation}
Including the source-backed longitudinal tip field gives about 0.263. These numbers are orientation estimates rather than a fit to the non-equilibrium STM electron populations.

The same experiment cites an electron lifetime of order 100 ns on the atom, while $\fR^{99}=21.93$ kHz corresponds to $t_g=11.4\,\mu$s. If the static-manifold propagation is to be interpreted as a gate without switching, an exponential switching time $\tau_e$ must satisfy
\begin{equation}
1-e^{-t_g/\tau_e}\ll1.
\end{equation}
Numerically, $P_{\rm switch}<0.1$ requires $\tau_e>108\,\mu$s, and $P_{\rm switch}<0.01$ requires $\tau_e>1.14$ ms. These are control specifications, not demonstrated electron lifetimes. The pulsed readout protocol demonstrates that DC and RF excitation can be turned off between probe pulses \cite{stolte2025}; whether a control configuration can also suppress electron switching to the required scale remains an experimental question.

If the electron is prepared and remains in one manifold, the waveform contains seven rather than fourteen resonant tones and the inter-manifold ceiling $S/2$ no longer binds. The leading-order spacing relevant to adjacent transitions is then $S$. No seven-tone gate-rate threshold is inferred here because it requires a dedicated propagation.

\section{Coherence, linewidth and next measurements}
No measured \Ti\ nuclear $T_2$ is assumed. Under the adjacent-level pure-dephasing model of the archived analysis, $T_{2,{\rm adj}}/t_g$ is approximately $10^3$, $10^4$ and $10^5$ for 99\%, 99.9\% and 99.99\% target fidelities. At 20 kHz, these correspond to 12.6 ms, 126 ms and 1.26 s.

A future selectivity figure of merit is $Q_{\rm sel}=\DeltaInter/\Gamma_{\rm NMR}$, where $\Gamma_{\rm NMR}$ is a linewidth measured under the relevant drive and junction conditions. The fit precision reported in the $^{47}$Ti single-atom NMR preprint is not used as a linewidth \cite{vennema2025}. The highest-leverage measurements are therefore: (i) direct \Ti\ NMR through 0.99 to 1.07 T; (ii) direct \Ti\ quadrupole and transverse-hyperfine extraction on the target junction; (iii) time-domain nuclear Rabi calibration; (iv) Ramsey and echo $T_2$; (v) electron switching rates during the control interval; and (vi) adsorption-site or tip-configuration dependence of the EFG.

\section{Reproducibility package}
The version 2.0.0 repository-ready folder accompanying this revision contains the low-field 16-level scan, validation records, dense control sweeps, one-at-a-time sensitivity data, field-quadrupole design map, scripts, publication figures and SHA-256 checksums. The earlier Zenodo release remains available as version 1.0.0 at DOI \href{https://doi.org/10.5281/zenodo.21969250}{10.5281/zenodo.21969250}; the version 2.0.0 release should be created under the same concept record before final publication.

\begin{thebibliography}{99}
\bibitem{stolte2025} Stolte E W, Lee J, Vennema H G, Broekhoven R, Teng E, Katan A J, Veldman L M, Willke P and Otte S 2025 \emph{Nat. Commun.} \textbf{16} 7785 doi:10.1038/s41467-025-63232-5
\bibitem{vennema2025} Vennema H G, Mier C, Stolte E W, Edens L, Lee J and Otte S 2025 Nuclear magnetic resonance on a single atom with a local probe arXiv:2512.11652
\bibitem{kim2022} Kim J, Noh K, Chen Y, Donati F, Heinrich A J, Wolf C and Bae Y 2022 \emph{Nano Lett.} \textbf{22} 9766--9772 doi:10.1021/acs.nanolett.2c02782
\bibitem{stone2016} Stone N J 2016 \emph{At. Data Nucl. Data Tables} \textbf{111--112} 1--28 doi:10.1016/j.adt.2015.12.002
\bibitem{asaad2020} Asaad S \emph{et al.} 2020 \emph{Nature} \textbf{579} 205--209 doi:10.1038/s41586-020-2057-7
\bibitem{fernandez2024} Fern\'andez de Fuentes I \emph{et al.} 2024 \emph{Nat. Commun.} \textbf{15} 1380 doi:10.1038/s41467-024-45368-y
\bibitem{yu2025} Yu X \emph{et al.} 2025 \emph{Nat. Phys.} \textbf{21} 362--367 doi:10.1038/s41567-024-02745-0
\bibitem{nielsen2002} Nielsen M A 2002 \emph{Phys. Lett. A} \textbf{303} 249--252 doi:10.1016/S0375-9601(02)01272-0
\end{thebibliography}
\end{document}
